\documentclass[10pt,a4paper]{article}

\usepackage[utf8]{inputenc}
\usepackage[T1]{fontenc}

\usepackage{amsmath,mathtools,amsfonts,amssymb}
\usepackage{graphicx}
\usepackage[spanish]{babel}
\usepackage{lastpage,tabto,xcolor}
\usepackage[a4paper, top=0.8in, bottom=1in, left=0.7in, right=0.7in]{geometry}
\usepackage{fancyhdr}
\usepackage{comment}
\usepackage{multicol}
\usepackage{titling} % Para ajustar los espacios del título
\usepackage{tasks} % Para listas en columnas
\usepackage{float}
\usepackage[mode=image]{standalone}
\usepackage{tikz}
\usepackage{pgfplots}
\pgfplotsset{width=7cm,compat=1.15}
\usepackage{caption}

% Fecha de versión automática según fecha de compilación
\newcommand{\verdate}{\the\year.\ifnum\month<10 0\fi\the\month.\ifnum\day<10 0\fi\the\day}

% Ajusta la separación entre el encabezado y el contenido
\setlength{\headsep}{10pt}

% Ajuste del título:
\setlength{\droptitle}{-0.5in} % Mueve el título ligeramente hacia arriba
\pretitle{\begin{center}\vspace{-1em}\normalfont\Large}
	\posttitle{\par\end{center}\vspace{-4em}}
\preauthor{\begin{center}\vspace{-0.5em}}
	\postauthor{\par\end{center}}
\predate{}\postdate{}

\pagestyle{fancy}
\lhead{\footnotesize Análisis de Señales y Sistemas  - Año \the\year{} Ver. \verdate}
\rhead{\footnotesize Trabajo Práctico N$^\text{o}$7: Sistemas de tiempo discreto.}
\rfoot{\footnotesize Página \thepage\ de \pageref{LastPage}}
\renewcommand{\headrulewidth}{0.4pt}
\renewcommand{\footrulewidth}{0.4pt}

% Título optimizado con espacios reducidos
\title{%
	{\normalsize Departamento de Ingeniería en Tecnologías Electrónicas, UTN FRM}\\[0.1cm]
	{\large Análisis de Señales y Sistemas}\\[0.15cm]
	{\normalsize Trabajo Práctico N$^{\text o}$7: Sistemas de tiempo discreto}%
}
\author{}
\date{}


\begin{document}
	\maketitle
	\thispagestyle{fancy}

	% Entorno de dos columnas
	\begin{multicols}{2}

		\label{sec:enun}

		{\small\noindent Salvo indicación en contrario, los sistemas
		descritos por una ecuación en diferencias se toman en reposo
		inicial.}

		\vspace{0.4em}

		% ***********************************************
		% 		Propiedades de los sistemas discretos
		% ***********************************************
		\textbf{Propiedades de los sistemas discretos}
		\begin{enumerate}

		% ***********************************************
		% 				Ejercicio 1
		% ***********************************************
		\item Sean los sistemas
		\begin{align*}
			\mathcal{T}_1&: \; y[n] = x[n] - 2\,x[n-2], \\
			\mathcal{T}_2&: \; y[n] = x[n]\,x[n-1], \\
			\mathcal{T}_3&: \; y[n] = \textstyle\sum_{k=n-2}^{n+2} x[k], \\
			\mathcal{T}_4&: \; y[n] = (n+1)\,x[n].
		\end{align*}
		\begin{enumerate}
			\item Indicar cuáles tienen memoria y cuáles son causales.
			\item Determinar cuáles son lineales y cuáles invariantes en el
			tiempo, exhibiendo un contraejemplo en cada caso negativo.
			\item Determinar cuáles son estables BIBO.
			\item Señalar cuál de los cuatro no puede usarse en un
			procesamiento en tiempo real e indicar en qué situación sí es
			utilizable.
		\end{enumerate}

		% ***********************************************
		% 		Respuesta al impulso y convolución
		% ***********************************************
		\textbf{Respuesta al impulso y convolución}

		% ***********************************************
		% 				Ejercicio 2
		% ***********************************************
		\item Un sistema LIT responde a la ecuación
		\[
			y[n] = 2\,x[n] + x[n-1] - x[n-3],
		\]
		y se lo excita con la entrada
		\[
			x[n] = \{\, \underset{\uparrow}{1} \quad 2 \quad -1 \,\}.
		\]
		\begin{enumerate}
			\item Obtener y graficar $h[n]$, indicando su soporte.
			\item Calcular $y[n]$ armando la tabla de copias desplazadas de
			$h[n]$, con una fila por muestra no nula de la entrada.
			\item Recalcular la salida completa con el procedimiento de
			invertir y desplazar, graficando $x[k]$ y $h[n-k]$ en cada
			instante, y verificar que coincide con la última fila de la
			tabla.
			\item Contrastar el largo del soporte de la salida con el que
			predice la regla $L_x + L_h - 1$.
			\item Rehacer el cálculo como producto de polinomios y comparar
			los coeficientes con la última fila de la tabla.
		\end{enumerate}

		% ***********************************************
		% 				Ejercicio 3
		% ***********************************************
		\item Sean las convoluciones
		\begin{align*}
			y_1[n] &= \bigl(u[n]-u[n-4]\bigr) * \bigl(u[n]-u[n-3]\bigr), \\
			y_2[n] &= \left(\tfrac{1}{2}\right)^{\!n} u[n] *
			          \left(\delta[n]-\tfrac{1}{2}\,\delta[n-1]\right), \\
			y_3[n] &= \left(\tfrac{1}{2}\right)^{\!n} u[n] *
			          \left(\tfrac{1}{4}\right)^{\!n} u[n].
		\end{align*}
		\begin{enumerate}
			\item Calcular y graficar $y_1[n]$, indicando su soporte.
			\item Calcular $y_2[n]$ e indicar qué relación guardan entre sí
			los dos sistemas cuyas respuestas al impulso se convolucionaron.
			\item Calcular $y_3[n]$ en forma cerrada e indicar cuál de las dos
			exponenciales gobierna el decaimiento cuando $n$ es grande.
		\end{enumerate}

		% ***********************************************
		% 				Ejercicio 4
		% ***********************************************
		\item Un sistema LIT tiene respuesta al impulso
		$h[n]=\left(\tfrac{1}{2}\right)^{n} u[n]$. Se llama $g[n]$ a la salida
		que produce ante la entrada $u[n]$.
		\begin{enumerate}
			\item Resolver la suma de convolución y escribir $g[n]$ en forma
			cerrada.
			\item Mostrar, sobre esa misma suma, que $g[n]$ es la suma
			acumulada de los valores $h[k]$ con $k \le n$.
			\item Verificar que $h[n]=g[n]-g[n-1]$ y explicar por qué es la
			primera diferencia la que recupera $h[n]$.
			\item Determinar hacia qué valor tiende $g[n]$ al aumentar $n$ y
			relacionar ese límite con la estabilidad del sistema.
		\end{enumerate}

		% ***********************************************
		% 		Interconexión de sistemas LIT
		% ***********************************************
		\textbf{Interconexión de sistemas LIT}

		% ***********************************************
		% 				Ejercicio 5
		% ***********************************************
		\item El sistema de la Figura~\ref{fig:tp7_interconexion} se arma con
		tres sistemas LIT de respuestas al impulso
		\begin{align*}
			h_1[n] &= \delta[n]+\delta[n-1], \\
			h_2[n] &= \delta[n]-\delta[n-1], \\
			h_3[n] &= \left(\tfrac{1}{2}\right)^{\!n} u[n].
		\end{align*}

		\begin{center}
			\includestandalone[width=0.98\linewidth]{figures/fig_tp7_interconexion/fig_tp7_interconexion}
			\captionof{figure}{Interconexión del ejercicio 5.}
			\label{fig:tp7_interconexion}
		\end{center}

		\begin{enumerate}
			\item Hallar la respuesta al impulso equivalente de las dos ramas
			en paralelo e indicar qué operación aplica esa etapa sobre la
			entrada.
			\item Hallar la respuesta al impulso $h[n]$ del sistema completo.
			\item Clasificar el sistema completo como FIR o IIR y determinar
			si es estable.
			\item Indicar si $h[n]$ cambia al intercambiar el orden de las dos
			etapas de la cascada. Justificar.
		\end{enumerate}

		% ***********************************************
		% 		Propiedades leídas en h[n]
		% ***********************************************
		\textbf{Propiedades leídas en $h[n]$}

		% ***********************************************
		% 				Ejercicio 6
		% ***********************************************
		\item Sean las respuestas al impulso
		\begin{align*}
			h_1[n] &= 3\,\delta[n], &
			h_2[n] &= u[n]-u[n-5], \\
			h_3[n] &= \left(\tfrac{3}{2}\right)^{\!n} u[n], &
			h_4[n] &= 2^{n}\,u[-n], \\
			h_5[n] &= \cos\!\left(\tfrac{\pi}{4}\,n\right) u[n]. &
		\end{align*}
		\begin{enumerate}
			\item Indicar cuáles corresponden a sistemas sin memoria y cuáles
			a sistemas causales.
			\item Determinar cuáles corresponden a sistemas estables BIBO.
			\item Clasificar cada sistema como FIR o IIR.
			\item Señalar el sistema que resulta estable sin ser causal y
			explicar por qué las dos condiciones son independientes entre sí.
		\end{enumerate}

		% ***********************************************
		% 		Ecuaciones en diferencias y realización
		% ***********************************************
		\textbf{Ecuaciones en diferencias y realización}

		% ***********************************************
		% 				Ejercicio 7
		% ***********************************************
		\item Un sistema en reposo inicial ($y[-1]=0$) responde a la
		ecuación en diferencias
		\[
			y[n] = x[n] + 0{,}8\,y[n-1].
		\]
		\begin{enumerate}
			\item Correr la recursión con $x[n]=\delta[n]$ y escribir $h[n]$
			en forma cerrada.
			\item Clasificar el sistema como FIR o IIR y determinar si es
			estable.
			\item Calcular las primeras cinco muestras de la salida ante
			$x[n]=u[n]$ y determinar hacia qué valor tiende.
			\item Graficar la realización del sistema con retardos unitarios,
			sumadores y escaladores.
			\item Determinar qué valores del coeficiente que multiplica a
			$y[n-1]$ producen un sistema inestable.
		\end{enumerate}

		% ***********************************************
		% 				Ejercicio 8
		% ***********************************************
		\item Un sistema en reposo inicial ($y[-1]=y[-2]=0$) responde a la
		ecuación en diferencias
		\[
			y[n] = x[n] + 0{,}9\,y[n-1] - 0{,}2\,y[n-2].
		\]
		\begin{enumerate}
			\item Correr la recursión con $x[n]=\delta[n]$ y escribir las
			primeras seis muestras de $h[n]$.
			\item Clasificar el sistema como FIR o IIR e indicar cuántas
			condiciones iniciales necesita la recursión para arrancar.
			\item Calcular las primeras cinco muestras de la salida ante
			$x[n]=u[n]$ y determinar hacia qué valor tiende, sin resolver la
			recursión en forma cerrada.
			\item Graficar la realización del sistema con retardos unitarios,
			sumadores y escaladores.
			\item Rehacer las primeras seis muestras de $h[n]$ con el
			coeficiente de $y[n-2]$ cambiado de $-0{,}2$ a $+0{,}2$, indicar
			qué le pasa a la salida y qué responde en ese caso el
			procedimiento del inciso c) sobre el valor al que tiende.
		\end{enumerate}

		% ***********************************************
		% 				Ejercicio 9
		% ***********************************************
		\item Sean los sistemas
		\begin{align*}
			\mathcal{T}_A&: \; y[n] = x[n]+x[n-1]+x[n-2]+x[n-3], \\
			\mathcal{T}_B&: \; y[n] = y[n-1] + x[n] - x[n-4].
		\end{align*}
		\begin{enumerate}
			\item Obtener $h[n]$ de cada uno corriendo la recursión con
			$x[n]=\delta[n]$, tomando $y[-1]=0$ en $\mathcal{T}_B$.
			\item Indicar si los dos sistemas son el mismo y justificar la
			respuesta sobre las dos respuestas al impulso.
			\item Graficar la realización de cada ecuación e indicar cuántos
			retardos y cuántas sumas usa cada una.
			\item Generalizar las dos ecuaciones a una ventana de $L$
			muestras, contar las sumas que cada forma hace por muestra de
			salida e indicar cuál conviene implementar cuando $L$ es grande.
		\end{enumerate}

		\end{enumerate}

	\end{multicols}

\end{document}
